
\documentclass[aspectratio=169,8pt]{beamer}
\usetheme{Madrid}
\usepackage[utf8]{inputenc}
\usepackage[T1]{fontenc}
\usepackage{amsmath,amssymb}
\usepackage{booktabs}
\usepackage{enumitem}
\setlist[enumerate]{label=\Alph*.,leftmargin=*,itemsep=1pt,topsep=2pt}
\setlist[itemize]{leftmargin=*,itemsep=1pt,topsep=2pt}
\setbeamerfont{block body}{size=\tiny}
\setbeamerfont{block title}{size=\scriptsize}
\setbeamerfont{frametitle}{size=\footnotesize}
\setbeamerfont{normal text}{size=\tiny}
\setbeamertemplate{navigation symbols}{}
\setbeamertemplate{itemize item}{\tiny$\bullet$}
\setbeamertemplate{itemize subitem}{\tiny$\circ$}

\title[GARP RAI]{GARP Risk \& AI --- Q351--Q450}
\subtitle{Unofficial Exam Prep}
\author{Unofficial}
\date{\today}

\begin{document}
	
	\begin{frame}
		\titlepage
	\end{frame}
	
	% =================================================================
	\section{Data Governance (continued)}
	% =================================================================
	



\begin{frame}{Q520 --- Demographic Parity}
	\begin{block}{Question}
		Demographic parity requires equality of:
		\begin{enumerate}
			\item True positive rates
			\item False positive rates
			\item Positive prediction rates
			\item Precision
		\end{enumerate}
	\end{block}
	\begin{exampleblock}{Answer: C}\end{exampleblock}
	\begin{block}{Explanation}
		Demographic parity requires:
		\[
		P(\hat{Y}=1|A=a)
		=
		P(\hat{Y}=1|A=b)
		\]
		Positive prediction rates must be equal across groups.
	\end{block}
\end{frame}

\begin{frame}{Q521 --- Equal Opportunity}
	\begin{block}{Question}
		Equal opportunity requires equality of:
		\begin{enumerate}
			\item Precision
			\item Sensitivity (True Positive Rate)
			\item Accuracy
			\item Positive prediction rate
		\end{enumerate}
	\end{block}
	\begin{exampleblock}{Answer: B}\end{exampleblock}
	\begin{block}{Explanation}
		Equal opportunity requires:
		\[
		P(\hat{Y}=1|Y=1,A)
		\]
		to be equal across groups.
	\end{block}
\end{frame}

\begin{frame}{Q522 --- Predictive Rate Parity}
	\begin{block}{Question}
		Predictive rate parity requires equality of:
		\begin{enumerate}
			\item Precision
			\item Recall
			\item False Positive Rate
			\item Accuracy
		\end{enumerate}
	\end{block}
	\begin{exampleblock}{Answer: A}\end{exampleblock}
	\begin{block}{Explanation}
		Predictive rate parity requires:
		\[
		P(Y=1|\hat{Y}=1,A)
		\]
		to be equal across groups.
	\end{block}
\end{frame}

\begin{frame}{Q523 --- Equalized Odds}
	\begin{block}{Question}
		Equalized odds requires equality of:
		\begin{enumerate}
			\item Precision and Accuracy
			\item TPR and FPR
			\item Precision and Recall
			\item Positive Prediction Rate and Precision
		\end{enumerate}
	\end{block}
	\begin{exampleblock}{Answer: B}\end{exampleblock}
	\begin{block}{Explanation}
		Equalized odds requires:
		\[
		TPR_A = TPR_B
		\]
		and
		\[
		FPR_A = FPR_B
		\]
		simultaneously.
	\end{block}
\end{frame}

\begin{frame}{Q524 --- Fairness Conditioning Trap}
	\begin{block}{Question}
		Which fairness metric conditions on predicted positives?
		\begin{enumerate}
			\item Demographic Parity
			\item Equal Opportunity
			\item Predictive Rate Parity
			\item Equalized Odds
		\end{enumerate}
	\end{block}
	\begin{exampleblock}{Answer: C}\end{exampleblock}
	\begin{block}{Explanation}
		Predictive rate parity evaluates:
		\[
		P(Y=1|\hat{Y}=1)
		\]
		which conditions on predicted positives.
	\end{block}
\end{frame}

\begin{frame}{Q525 --- Precision}
	\begin{block}{Question}
		Precision answers which question?
		\begin{enumerate}
			\item Of all actual positives, how many were found?
			\item Of all predicted positives, how many were correct?
			\item Of all negatives, how many were correct?
			\item Of all predictions, how many were correct?
		\end{enumerate}
	\end{block}
	\begin{exampleblock}{Answer: B}\end{exampleblock}
	\begin{block}{Explanation}
		\[
		Precision = \frac{TP}{TP+FP}
		\]
	\end{block}
\end{frame}

\begin{frame}{Q526 --- Recall / Sensitivity}
	\begin{block}{Question}
		Recall (Sensitivity) answers which question?
		\begin{enumerate}
			\item Of all actual positives, how many were correctly identified?
			\item Of all predicted positives, how many were correct?
			\item Of all negatives, how many were correct?
			\item Of all predictions, how many were correct?
		\end{enumerate}
	\end{block}
	\begin{exampleblock}{Answer: A}\end{exampleblock}
	\begin{block}{Explanation}
		\[
		Recall = \frac{TP}{TP+FN}
		\]
	\end{block}
\end{frame}

\begin{frame}{Q527 --- False Positive Rate}
	\begin{block}{Question}
		False Positive Rate (FPR) measures:
		\begin{enumerate}
			\item The fraction of actual positives missed
			\item The fraction of actual negatives incorrectly predicted positive
			\item The fraction of predicted positives that are wrong
			\item The fraction of predictions that are correct
		\end{enumerate}
	\end{block}
	\begin{exampleblock}{Answer: B}\end{exampleblock}
	\begin{block}{Explanation}
		\[
		FPR=\frac{FP}{FP+TN}
		\]
	\end{block}
\end{frame}

\begin{frame}{Q528 --- ROC Curve}
	\begin{block}{Question}
		An ROC curve plots:
		\begin{enumerate}
			\item Precision vs Recall
			\item Accuracy vs Threshold
			\item TPR vs FPR
			\item FPR vs Precision
		\end{enumerate}
	\end{block}
	\begin{exampleblock}{Answer: C}\end{exampleblock}
	\begin{block}{Explanation}
		The ROC curve plots:
		\[
		TPR
		\]
		against
		\[
		FPR
		\]
		across thresholds.
	\end{block}
\end{frame}

\begin{frame}{Q529 --- AUC}
	\begin{block}{Question}
		AUC is BEST interpreted as:
		\begin{enumerate}
			\item Average precision
			\item Average accuracy
			\item Overall ability to distinguish positives from negatives
			\item Probability that a prediction is correct
		\end{enumerate}
	\end{block}
	\begin{exampleblock}{Answer: C}\end{exampleblock}
\end{frame}

\begin{frame}{Q530 --- Perfect AUC}
	\begin{block}{Question}
		What AUC corresponds to a perfect classifier?
		\begin{enumerate}
			\item 0
			\item 0.5
			\item 0.75
			\item 1.0
		\end{enumerate}
	\end{block}
	\begin{exampleblock}{Answer: D}\end{exampleblock}
\end{frame}

\begin{frame}{Q531 --- Random Guessing}
	\begin{block}{Question}
		An AUC of 0.5 suggests:
		\begin{enumerate}
			\item Perfect classification
			\item Random guessing
			\item Severe overfitting
			\item High precision
		\end{enumerate}
	\end{block}
	\begin{exampleblock}{Answer: B}\end{exampleblock}
\end{frame}

\begin{frame}{Q532 --- Threshold Effect}
	\begin{block}{Question}
		Lowering a decision threshold will generally:
		\begin{enumerate}
			\item Increase Recall and Increase FPR
			\item Increase Recall and Decrease FPR
			\item Decrease Recall and Increase FPR
			\item Decrease Recall and Decrease FPR
		\end{enumerate}
	\end{block}
	\begin{exampleblock}{Answer: A}\end{exampleblock}
	\begin{block}{Explanation}
		A lower threshold classifies more observations as positive, increasing both true positives and false positives.
	\end{block}
\end{frame}

\begin{frame}{Q533 --- Threshold Tradeoff}
	\begin{block}{Question}
		A fraud team lowers its model threshold. What is the MOST likely outcome?
		\begin{enumerate}
			\item More fraud caught and more false alarms
			\item Less fraud caught and fewer false alarms
			\item Higher precision and lower recall
			\item Lower recall and higher specificity
		\end{enumerate}
	\end{block}
	\begin{exampleblock}{Answer: A}\end{exampleblock}
\end{frame}





\begin{frame}{Q510 --- Predictive Rate Parity}
	\begin{block}{Question}
		Predictive rate parity requires equality of:
		\begin{enumerate}
			\item Positive prediction rates
			\item Precision across groups
			\item Sensitivity across groups
			\item False positive rates across groups
		\end{enumerate}
	\end{block}
	\begin{exampleblock}{Answer: B}
		Precision across groups
	\end{exampleblock}
	\begin{block}{Explanation}
		Predictive rate parity (also called \textbf{conditional use accuracy equality} or
		\textbf{predictive parity}) requires that the \textbf{positive predictive value (precision)}
		be equal across groups:
		\[
		P(Y = 1 \mid \hat{Y} = 1, A = 0) = P(Y = 1 \mid \hat{Y} = 1, A = 1)
		\]
		In words: among all individuals who are \emph{predicted} to be positive,
		the proportion who are \emph{actually} positive must be the same for every group.
		This conditions on the \textbf{prediction} (the predicted positive class), not on the
		true label. It does \emph{not} require equal positive prediction rates
		(that is demographic parity), nor equal sensitivity (that is equal opportunity),
		nor equal false positive rates (that is predictive equality).
	\end{block}
\end{frame}

\begin{frame}{Q511 --- Equal Opportunity}
	\begin{block}{Question}
		Equal opportunity requires equality of:
		\begin{enumerate}
			\item Precision
			\item Overall accuracy
			\item Sensitivity (true positive rate)
			\item Positive prediction rate
		\end{enumerate}
	\end{block}
	\begin{exampleblock}{Answer: C}
		Sensitivity (true positive rate)
	\end{exampleblock}
	\begin{block}{Explanation}
		Equal opportunity requires that the \textbf{true positive rate (sensitivity / recall)}
		be equal across groups:
		\[
		P(\hat{Y} = 1 \mid Y = 1, A = 0) = P(\hat{Y} = 1 \mid Y = 1, A = 1)
		\]
		In words: among all individuals who are \emph{truly} positive,
		the proportion who are \emph{correctly} identified as positive must be the same
		for every group. This conditions on the \textbf{true label} (the actual positive class).
		It does \emph{not} require equal precision (that is predictive rate parity),
		nor equal overall accuracy, nor equal positive prediction rates
		(that is demographic parity).
	\end{block}
\end{frame}

\begin{frame}{Q512 --- Direction of Conditioning}
	\begin{block}{Question}
		Which fairness criterion conditions on the predicted positive class?
		\begin{enumerate}
			\item Equal Opportunity
			\item Predictive Rate Parity
			\item Equalized Odds
			\item Demographic Parity
		\end{enumerate}
	\end{block}
	\begin{exampleblock}{Answer: B}
		Predictive Rate Parity
	\end{exampleblock}
	\begin{block}{Explanation}
		Predictive rate parity conditions on the \textbf{predicted positive class}
		(\(\hat{Y} = 1\)). It asks: given that the model predicted positive,
		what is the probability that the individual is truly positive?
		\[
		P(Y = 1 \mid \hat{Y} = 1, A = 0) = P(Y = 1 \mid \hat{Y} = 1, A = 1)
		\]
		This is exactly the definition of \textbf{precision} (positive predictive value).
		By contrast:
		\begin{itemize}
			\item \textbf{Equal Opportunity} conditions on the \emph{true} positive class
			(\(Y = 1\)) and equalizes sensitivity (TPR).
			\item \textbf{Equalized Odds} conditions on \emph{both} the true positive and
			true negative classes, equalizing TPR and FPR simultaneously.
			\item \textbf{Demographic Parity} does not condition on any outcome label;
			it simply equalizes the positive prediction rate \(P(\hat{Y} = 1 \mid A)\).
		\end{itemize}
	\end{block}
\end{frame}

\begin{frame}{Q513 --- Reverse Conditioning Trap}
	\begin{block}{Question}
		Which fairness criterion conditions on those who are truly positive?
		\begin{enumerate}
			\item Equal Opportunity
			\item Predictive Rate Parity
			\item Demographic Parity
			\item Predictive Equality
		\end{enumerate}
	\end{block}
	\begin{exampleblock}{Answer: A}
		Equal Opportunity
	\end{exampleblock}
	\begin{block}{Explanation}
		Equal opportunity conditions on the \textbf{true positive class} (\(Y = 1\)).
		It requires equal sensitivity (true positive rate) across groups:
		\[
		P(\hat{Y} = 1 \mid Y = 1, A = 0) = P(\hat{Y} = 1 \mid Y = 1, A = 1)
		\]
		In words: among those who are \emph{truly} positive, the probability of being
		correctly predicted positive must be the same for all groups.
		
		\textbf{Why the others are wrong:}
		\begin{itemize}
			\item \textbf{Predictive Rate Parity} conditions on the \emph{predicted}
			positive class (\(\hat{Y} = 1\)), not the true positive class.
			\item \textbf{Demographic Parity} does not condition on any outcome;
			it equalizes the positive prediction rate.
			\item \textbf{Predictive Equality} conditions on the \emph{true negative}
			class (\(Y = 0\)) and equalizes the false positive rate.
		\end{itemize}
		This question is a ``reverse conditioning trap'' because it swaps the
		conditioning direction relative to Q512.
	\end{block}
\end{frame}

\begin{frame}{Q514 --- Demographic Parity}
	\begin{block}{Question}
		Demographic parity equalizes:
		\begin{enumerate}
			\item Precision
			\item Sensitivity
			\item Positive prediction rates
			\item Accuracy
		\end{enumerate}
	\end{block}
	\begin{exampleblock}{Answer: C}
		Positive prediction rates
	\end{exampleblock}
	\begin{block}{Explanation}
		Demographic parity (also called \textbf{statistical parity} or
		\textbf{independence}) requires that the \textbf{positive prediction rate}
		be equal across groups:
		\[
		P(\hat{Y} = 1 \mid A = 0) = P(\hat{Y} = 1 \mid A = 1)
		\]
		In words: the proportion of individuals predicted positive must be the same
		for every group, regardless of the true labels.
		
		\textbf{Why the others are wrong:}
		\begin{itemize}
			\item \textbf{Precision} is equalized by predictive rate parity.
			\item \textbf{Sensitivity} (TPR) is equalized by equal opportunity.
			\item \textbf{Accuracy} is not the target of any of these standard
			group-fairness criteria (though it may be equal in some settings).
		\end{itemize}
		Demographic parity is the only criterion here that can be evaluated
		\emph{without} observing ground-truth outcomes, because it depends only on
		predictions and group membership.
	\end{block}
\end{frame}

\begin{frame}{Q515 --- Fairness Mapping}
	\begin{block}{Question}
		Which pairing is correct?
		\begin{enumerate}
			\item Demographic Parity -- Equal Precision
			\item Equal Opportunity -- Equal Sensitivity
			\item Predictive Rate Parity -- Equal Recall
			\item Predictive Equality -- Equal Precision
		\end{enumerate}
	\end{block}
	\begin{exampleblock}{Answer: B}
		Equal Opportunity -- Equal Sensitivity
	\end{exampleblock}
	\begin{block}{Explanation}
		The correct pairing is \textbf{Equal Opportunity -- Equal Sensitivity}.
		Equal opportunity requires equal true positive rates (sensitivity/recall)
		across groups.
		
		\textbf{Why the others are wrong:}
		\begin{itemize}
			\item \textbf{A:} Demographic parity equalizes \emph{positive prediction rates},
			not precision.
			\item \textbf{C:} Predictive rate parity equalizes \emph{precision}
			(positive predictive value), not recall (sensitivity).
			\item \textbf{D:} Predictive equality equalizes \emph{false positive rates},
			not precision.
		\end{itemize}
		A useful memory aid:
		\begin{center}
			\begin{tabular}{ll}
				\textbf{Criterion} & \textbf{Equalized Metric} \\
				\hline
				Demographic Parity & Positive Prediction Rate \\
				Equal Opportunity & True Positive Rate (Sensitivity/Recall) \\
				Predictive Rate Parity & Precision (PPV) \\
				Predictive Equality & False Positive Rate \\
				Equalized Odds & TPR \emph{and} FPR \\
			\end{tabular}
		\end{center}
	\end{block}
\end{frame}

\begin{frame}{Q516 --- Predictive Equality}
	\begin{block}{Question}
		Predictive equality requires equality of:
		\begin{enumerate}
			\item False positive rates
			\item Precision
			\item Accuracy
			\item Positive prediction rates
		\end{enumerate}
	\end{block}
	\begin{exampleblock}{Answer: A}
		False positive rates
	\end{exampleblock}
	\begin{block}{Explanation}
		Predictive equality requires that the \textbf{false positive rate (FPR)}
		be equal across groups:
		\[
		P(\hat{Y} = 1 \mid Y = 0, A = 0) = P(\hat{Y} = 1 \mid Y = 0, A = 1)
		\]
		In words: among all individuals who are \emph{truly negative},
		the proportion who are \emph{incorrectly} predicted positive must be the
		same for every group. This conditions on the \textbf{true negative class}
		(\(Y = 0\)).
		
		\textbf{Why the others are wrong:}
		\begin{itemize}
			\item \textbf{Precision} is equalized by predictive rate parity.
			\item \textbf{Accuracy} is not a group-fairness criterion in this family.
			\item \textbf{Positive prediction rates} are equalized by demographic parity.
		\end{itemize}
		Predictive equality is sometimes called \textbf{equal false positive rate}
		and is one half of equalized odds (the other half being equal opportunity).
	\end{block}
\end{frame}

\begin{frame}{Q517 --- Equalized Odds}
	\begin{block}{Question}
		Equalized odds requires equality of:
		\begin{enumerate}
			\item Precision and accuracy
			\item Positive prediction rates and precision
			\item True positive rates and false positive rates
			\item Accuracy and sensitivity
		\end{enumerate}
	\end{block}
	\begin{exampleblock}{Answer: C}
		True positive rates and false positive rates
	\end{exampleblock}
	\begin{block}{Explanation}
		Equalized odds requires that \textbf{both} the true positive rate (TPR)
		\textbf{and} the false positive rate (FPR) be equal across groups:
		\[
		P(\hat{Y} = 1 \mid Y = 1, A = 0) = P(\hat{Y} = 1 \mid Y = 1, A = 1)
		\]
		\[
		P(\hat{Y} = 1 \mid Y = 0, A = 0) = P(\hat{Y} = 1 \mid Y = 0, A = 1)
		\]
		In words: the model must have equal sensitivity \emph{and} equal
		false positive rate for every group. Equivalently, the prediction
		\(\hat{Y}\) must be independent of the protected attribute \(A\)
		\emph{conditional on the true label} \(Y\).
		
		\textbf{Why the others are wrong:}
		\begin{itemize}
			\item \textbf{A:} Precision and accuracy are not the targets of equalized odds.
			\item \textbf{B:} Positive prediction rates and precision are targets of
			demographic parity and predictive rate parity, respectively.
			\item \textbf{D:} Accuracy and sensitivity are not jointly equalized by
			equalized odds (only TPR and FPR are).
		\end{itemize}
		Equalized odds is a \textbf{combination} of equal opportunity (equal TPR)
		and predictive equality (equal FPR).
	\end{block}
\end{frame}

\begin{frame}{Q518 --- Label Requirement}
	\begin{block}{Question}
		Which fairness criterion can be evaluated without observing ground-truth outcomes?
		\begin{enumerate}
			\item Equal Opportunity
			\item Predictive Rate Parity
			\item Equalized Odds
			\item Demographic Parity
		\end{enumerate}
	\end{block}
	\begin{exampleblock}{Answer: D}
		Demographic Parity
	\end{exampleblock}
	\begin{block}{Explanation}
		Demographic parity can be evaluated \textbf{without observing ground-truth
			outcomes} because it depends only on the \textbf{predicted labels} and
		\textbf{group membership}:
		\[
		P(\hat{Y} = 1 \mid A = 0) = P(\hat{Y} = 1 \mid A = 1)
		\]
		No true labels \(Y\) are needed.
		
		\textbf{Why the others require ground truth:}
		\begin{itemize}
			\item \textbf{Equal Opportunity} requires \(Y = 1\) to compute sensitivity.
			\item \textbf{Predictive Rate Parity} requires \(Y = 1\) to compute precision.
			\item \textbf{Equalized Odds} requires both \(Y = 1\) and \(Y = 0\) to
			compute TPR and FPR.
		\end{itemize}
		This is why demographic parity is often attractive in settings where
		labels are unavailable, delayed, or expensive to obtain. However, it can
		be satisfied even when the model is inaccurate, which is a known limitation.
	\end{block}
\end{frame}

\begin{frame}{Q519 --- Memorization Trap}
	\begin{block}{Question}
		A candidate says:
		``Predictive rate parity equals the share of actual positives correctly identified.''
		
		Which metric are they accidentally describing?
		\begin{enumerate}
			\item Predictive Rate Parity
			\item Demographic Parity
			\item Equal Opportunity
			\item Calibration
		\end{enumerate}
	\end{block}
	\begin{exampleblock}{Answer: C}
		Equal Opportunity
	\end{exampleblock}
	\begin{block}{Explanation}
		The phrase ``the share of actual positives correctly identified'' is the
		definition of \textbf{sensitivity} (true positive rate / recall):
		\[
		\text{Sensitivity} = P(\hat{Y} = 1 \mid Y = 1)
		\]
		Equal opportunity requires \emph{equal} sensitivity across groups.
		So the candidate is accidentally describing \textbf{Equal Opportunity},
		not predictive rate parity.
		
		\textbf{What predictive rate parity actually is:}
		Predictive rate parity requires equal \textbf{precision} (positive predictive
		value) across groups:
		\[
		P(Y = 1 \mid \hat{Y} = 1, A = 0) = P(Y = 1 \mid \hat{Y} = 1, A = 1)
		\]
		Precision is ``the share of \emph{predicted} positives that are actually positive,''
		which is the reverse conditioning of sensitivity.
		
		\textbf{Why the others are wrong:}
		\begin{itemize}
			\item \textbf{A:} Predictive rate parity is precision, not sensitivity.
			\item \textbf{B:} Demographic parity is the positive prediction rate,
			\(P(\hat{Y} = 1 \mid A)\).
			\item \textbf{D:} Calibration refers to the agreement between predicted
			probabilities and observed frequencies, not this quantity.
		\end{itemize}
		\textbf{Key distinction:}
		\begin{center}
			\begin{tabular}{ll}
				\textbf{Metric} & \textbf{Conditioning} \\
				\hline
				Sensitivity (Equal Opportunity) & \(P(\hat{Y}=1 \mid Y=1)\) \\
				Precision (Predictive Rate Parity) & \(P(Y=1 \mid \hat{Y}=1)\) \\
			\end{tabular}
		\end{center}
		The candidate has confused the direction of conditioning.
	\end{block}
\end{frame}
\begin{frame}
	\href{https://www.youtube.com/watch?v=D035nSq7LaU}{Watch the video}
\end{frame}





\begin{frame}{Q500 --- Ground Truth Availability}
	\begin{block}{Question}
		A bank wants to assess fairness immediately after making lending decisions. Loan defaults typically occur years later. Which fairness criterion is MOST practical at decision time?
		\begin{enumerate}
			\item Predictive Parity
			\item Equal Opportunity
			\item Demographic Parity
			\item Equalized Odds
		\end{enumerate}
	\end{block}
	\begin{exampleblock}{Answer: C}\end{exampleblock}
	\begin{block}{Explanation}
		Demographic parity does not require knowledge of future outcomes. The other metrics depend on ground-truth labels that become available only later.
	\end{block}
\end{frame}

\begin{frame}{Q501 --- Base Rates}
	\begin{block}{Question}
		Why can demographic parity be problematic in medical diagnosis?
		\begin{enumerate}
			\item It requires too much training data.
			\item It ignores legitimate differences in disease prevalence across groups.
			\item It requires causal models.
			\item It maximizes false negatives.
		\end{enumerate}
	\end{block}
	\begin{exampleblock}{Answer: B}\end{exampleblock}
	\begin{block}{Explanation}
		If disease prevalence differs across groups, forcing equal positive prediction rates may reduce predictive usefulness and medical effectiveness.
	\end{block}
\end{frame}

\begin{frame}{Q502 --- Predictive Quality}
	\begin{block}{Question}
		Which statement about demographic parity is TRUE?
		\begin{enumerate}
			\item It guarantees equal error rates.
			\item It guarantees equal predictive value.
			\item It says nothing about whether predictions are correct.
			\item It guarantees equal opportunity.
		\end{enumerate}
	\end{block}
	\begin{exampleblock}{Answer: C}\end{exampleblock}
	\begin{block}{Explanation}
		Demographic parity focuses only on equal prediction rates across groups and does not evaluate prediction quality.
	\end{block}
\end{frame}

\begin{frame}{Q503 --- Impossibility Theorem}
	\begin{block}{Question}
		When base rates differ across groups, what does the fairness impossibility theorem imply?
		\begin{enumerate}
			\item Demographic parity, predictive parity, and equal opportunity can generally be satisfied simultaneously.
			\item One fairness definition becomes objectively correct.
			\item Certain fairness criteria cannot all be satisfied at the same time.
			\item Fairness measurement becomes impossible.
		\end{enumerate}
	\end{block}
	\begin{exampleblock}{Answer: C}\end{exampleblock}
	\begin{block}{Explanation}
		The theorem highlights that competing fairness definitions may conflict when underlying group base rates differ.
	\end{block}
\end{frame}

\begin{frame}{Q504 --- Fraud Detection}
	\begin{block}{Question}
		A fraud model is deployed in an environment where confirmed fraud labels may take months to establish and many fraud cases are never confirmed. Why might demographic parity still be considered?
		\begin{enumerate}
			\item It guarantees fairness.
			\item It does not require reliable ground-truth labels.
			\item It guarantees equal precision.
			\item It eliminates bias.
		\end{enumerate}
	\end{block}
	\begin{exampleblock}{Answer: B}\end{exampleblock}
	\begin{block}{Explanation}
		Demographic parity can be evaluated without waiting for true fraud outcomes, making it attractive in settings where labels are delayed, incomplete, or extremely rare.
	\end{block}
\end{frame}








\end{document}